\documentclass{situationist-preprint}

\title{Anxiety and the Disclosure of the Nothing}
\msid{PHE-16-a}
\msfamily{PHE-16}
\msdiscipline{Phenomenology}
\mstype{Principal manuscript}
\mskeywords{anxiety; the nothing; Heidegger; Carnap; disclosure; moods}

\begin{document}
\maketitle

\begin{abstract}
Heidegger claims that anxiety discloses the nothing. Carnap took the claim as a paradigm of metaphysical nonsense, produced by treating \enquote{nothing} as a name. We argue that anxiety discloses the nothing as the withdrawal of beings as a whole, that the disclosure is genuine, and that it does not reveal a further entity. Anxiety differs from fear in having no object among beings; what it presents is that beings as a whole, with the significance that ordinarily holds us among them, slip away. That withdrawal is a feature of one's situation, which the mood presents and which can be presented correctly or incorrectly. Carnap's objection succeeds against the reading on which \enquote{the nothing} names an entity, and fails against the reading on which it names the withdrawal; we argue that the second is Heidegger's.
\end{abstract}

\section{The problem}

In \emph{Being and Time}, anxiety is the fundamental attunement in which Dasein is brought before itself \citep[\S 40]{Heidegger1927}. Unlike fear, which is fear of a determinate threatening being, anxiety has no such object: that in the face of which one is anxious is being-in-the-world as such, and in anxiety the totality of involvements of things at hand collapses into insignificance. In the inaugural lecture of 1929, Heidegger states the result in a form that became notorious: anxiety reveals the nothing; in anxiety beings as a whole slip away, and \enquote{the nothing itself nihilates} \citep{Heidegger1929}.

\citet{Carnap1932} argued that such sentences violate logical syntax. \enquote{Nothing} is a quantifier phrase: \enquote{nothing is outside} says that there is no $x$ such that $x$ is outside. Heidegger uses it as a name and attaches predicates to it, and his sentences cannot be translated into a correct language. They express an attitude towards life, perhaps, but assert nothing. The question is what anxiety, if it discloses anything, discloses. We argue that it discloses the nothing as the withdrawal of beings as a whole, that the disclosure is genuine, and that it reveals no further entity.

\section{Framework}

\begin{definition}[Disclosure]\label{d:disc}
A mood \emph{discloses} a feature of one's situation if it presents the situation as having that feature, and can do so correctly or incorrectly.
\end{definition}

\begin{definition}[Withdrawal of beings as a whole]\label{d:withdr}
Beings as a whole \emph{withdraw} for a subject when the significance by which beings ordinarily solicit and hold her, as things to be used, pursued, avoided or attended to, lapses for all of them together, while they remain present.
\end{definition}

\begin{principle}[No determinate object]\label{p:obj}
Anxiety, unlike fear, is directed at no determinate being; nothing within the world is that in the face of which one is anxious.
\end{principle}

\begin{principle}[Felt slipping away]\label{p:slip}
In anxiety, beings are present but felt as slipping away: they no longer offer a hold, and this is felt as a feature of all of them together, not of any in particular.
\end{principle}

Principles~\ref{p:obj} and~\ref{p:slip} are Heidegger's descriptions, and they are recognizable: in anxiety one cannot say what one is anxious about, and the familiar world is uncanny, there but no longer one's own \citep[\S 40]{Heidegger1927}.

\section{The argument}

\begin{lemma}\label{l:withdr}
Anxiety presents beings as a whole as withdrawing.
\end{lemma}

\begin{proof}
By Principle~\ref{p:slip}, beings in anxiety remain present but are felt as offering no hold, and this is felt as a feature of all of them together. That is the withdrawal of Definition~\ref{d:withdr}. By Principle~\ref{p:obj}, the mood does not present any one being as the source of the withdrawal, so what it presents is the withdrawal of beings as a whole, not a property of a particular being.
\end{proof}

\begin{lemma}\label{l:genuine}
The disclosure is genuine: anxiety can present the withdrawal correctly.
\end{lemma}

\begin{proof}
The significance that beings have for a subject depends on the projects in terms of which she understands herself, and those projects rest on no ground beyond her own existence: none is required by the beings themselves. This is a feature of her situation, and it obtains whether or not she is anxious. When anxiety presents all significance as slipping away together, it presents her situation as one in which the hold that beings have on her is not secured by them. That presentation is correct when the dependence is as described, and it would be incorrect for a being whose purposes were fixed by its environment. By Definition~\ref{d:disc}, anxiety discloses the withdrawal.
\end{proof}

\begin{theorem}\label{t:main}
Anxiety discloses the nothing as the withdrawal of beings as a whole. The disclosure is genuine, and it does not reveal a further entity.
\end{theorem}

\begin{proof}
Call what anxiety discloses \enquote{the nothing}, as Heidegger does. By Lemma~\ref{l:withdr}, what it presents is the withdrawal of beings as a whole; by Lemma~\ref{l:genuine}, the presentation can be correct. The withdrawal is not a being: it is a lapse of the significance of all beings, which are themselves still present, and a lapse of significance is not an additional item alongside them. Hence the nothing so understood is not a further entity.
\end{proof}

The theorem gives a sense to Heidegger's notorious verb. That \enquote{the nothing nihilates} says, on this reading, that the withdrawal is not something the subject does, as she might stop attending to things, but something that happens to her situation as a whole; beings are disclosed in their withdrawal as beings, rather than as the equipment of her projects.

\section{Objections and replies}

\begin{objection}[Carnap]
\enquote{The nothing} is formed by nominalizing a quantifier. Any sentence that predicates something of it is either nonsense or, rewritten correctly, says that there is nothing that has the predicate. Either way, there is nothing that anxiety discloses \citep{Carnap1932}.
\end{objection}

\begin{reply}
The objection is correct against the reading on which \enquote{the nothing} names an object of which properties are predicated. Theorem~\ref{t:main} rejects that reading. \enquote{The nothing}, in our reconstruction, abbreviates a description of a feature of situations, the withdrawal of beings as a whole, which can be stated without nominalizing a quantifier: for every being, its significance for the subject lapses, and the subject is aware of this lapse as a whole. The sentence is syntactically in order, and it is either true or false of a given anxious subject. Heidegger anticipated the objection and replied that logic presupposes the nothing rather than founding it \citep{Heidegger1929}; we do not need that reply. \citet{Friedman2000} shows that the disagreement between Carnap and Heidegger concerned the status of logic, and our reconstruction is neutral on that question.
\end{reply}

\begin{objection}[Reification]
If anxiety discloses something, it discloses something that is. If the nothing is disclosed, it is, and so it is an entity after all.
\end{objection}

\begin{reply}
Disclosure is of features of a situation, and not every feature is an entity. A shortage of water is disclosed by thirst without being a further thing besides the water and the thirsty animal. The withdrawal of beings as a whole is in this respect like a shortage: it is real, it can be presented correctly or not, and it is not an addition to the inventory of beings.
\end{reply}

\begin{objection}[Mere feeling]
Anxiety is a physiological state with no cognitive content; its apparent disclosures are projections of distress onto the world.
\end{objection}

\begin{reply}
The objection would apply to every mood, and PHE-21 argues that moods are pre-reflective disclosures of one's situation, fallible but not arbitrary. Anxiety can be pathological, as in an anxiety disorder in which the sense of groundlessness is out of all proportion to the subject's situation, and then its disclosure is incorrect. The existence of incorrect cases presupposes a standard of correctness, which Lemma~\ref{l:genuine} supplies.
\end{reply}

\section{Conclusion}

Anxiety presents beings as a whole as slipping away and as offering no hold, and this presentation can be correct. Anxiety therefore discloses the nothing as the withdrawal of beings as a whole. Carnap's critique refutes the reification of the nothing, not the disclosure; what anxiety reveals is a feature of one's situation, not a further entity.

\begin{thebibliography}{9}\small
\bibitem[Carnap(1932)]{Carnap1932} Carnap, R. (1932). \"Uberwindung der Metaphysik durch logische Analyse der Sprache. \emph{Erkenntnis}, 2, 219--241.
\bibitem[Friedman(2000)]{Friedman2000} Friedman, M. (2000). \emph{A Parting of the Ways: Carnap, Cassirer, and Heidegger}. Open Court.
\bibitem[Heidegger(1927)]{Heidegger1927} Heidegger, M. (1927). \emph{Sein und Zeit}. Niemeyer. Trans. J. Macquarrie and E. Robinson, \emph{Being and Time}, Harper \& Row, 1962.
\bibitem[Heidegger(1929)]{Heidegger1929} Heidegger, M. (1929). \emph{Was ist Metaphysik?} Cohen.
\end{thebibliography}

\end{document}
